\section{Introduction}

% what we do and basic motivation for our asy framework
This paper considers inference in first-price and second-price sealed-bid auctions in empirical settings where we observe a possibly small number of auctions, each with a large number of bidders. The abundance of bidders in each auction prompts us to consider a novel asymptotic framework in which the number of bidders diverges, while allowing the number of auctions to be small and remain fixed. This framework differs substantially from the more conventional approach in which the researcher observes multiple bids from a large number of independent and identically distributed (i.i.d.) auctions. See \citet{athey2002,hailetamer2003,athey2007handbook,guerre2000}, among others. Since our asymptotic framework does not require the number of auctions to diverge, our analysis is suitable for applications with substantial heterogeneity across auctions, implying a limited number of truly homogeneous auctions. 

% what tests we have and why our method is good
Within our novel asymptotic framework, we introduce new inference methods for the winner's expected utility, the seller's expected revenue, and the valuation distribution's tail behavior. We show that the latter can be used to test the regularity conditions commonly assumed in auction literature. Our data requirements are minimal; our methods rely on observing transaction prices from a finite number of auctions. In particular, we do not need to observe multiple bids or the number of participating bidders in these auctions.

% empirical motivation
Our methodology is well-suited for auction settings where the number of participants is known to be large, and we only observe transaction prices for a possibly small number of homogeneous auctions. We illustrate our methodology using data from Hong Kong license plate auctions. Since 1973, the Hong Kong government has conducted monthly standard oral-ascending price auctions to sell vehicle license plates. Recent interviews with auction participants reveal that these auctions are well attended, with around 100 participants per auction. Only the transaction price is recorded.  Finally, license plates are inherently heterogeneous goods whose value depends on their configuration of letters and numbers. If we focus on license plates with specific configurations to get a homogeneous sample of auctions, the sample size is usually small. We use the Hong Kong vehicle license plate auction setup as our running example throughout the paper and our empirical application in Section \ref{sec:application}. Beyond this application, our econometric framework is also relevant for online auctions, treasury auctions, spectrum auctions, art auctions, and IPO auctions.
% Our motivating example is the Hong Kong vehicle license plate auction, conducted as a standard oral-ascending price auction.  
% For illustration, we use auctions for the extremely luxurious plates featuring only one letter as a running example. 
% Table \ref{tab:plate} in Section \ref{sec:setup} presents all four historical records of such auctions. 
% In these auctions, we only observe the transaction prices but not other information. The number of bidders is also unobserved but large (around 300 according to our interview of recent auction participants). 
% These features make econometric analysis highly nonstandard, rendering existing methods unsuitable, while our proposed method is well-suited. 
% Other relevant examples include online auctions, treasury auctions, spectrum auctions, art auctions, and IPO auctions, among others. 

% a bit of how the method works
%Our methodology characterizes the limiting behavior of transaction prices as the number of bidders increases. Unlike the traditional asymptotic framework with a growing number of auctions, the transaction price data does not allow us to fully identify the valuation distribution. However, it can reveal the distribution's tail properties, which are sufficient for conducting inference on practical objects of interest, such as the winner's expected utility and the seller's expected revenue. Moreover, the tail behavior allows us to evaluate whether the valuation distribution has bounded support and positive density at its highest value, which is a common regularity in the auction literature. \citep[e.g.,][]{MaskinRiley1984,guerre2000,GuerreLuo2022}. Within our asymptotic framework, our inference methods are shown to control size and exhibit desirable power properties. Specifically, our confidence intervals are shown to minimize expected length and our hypothesis tests are shown to maximize weighted average power. Finally, Monte Carlo simulations confirm the accuracy of our methodology in finite samples, and a study of Hong Kong car license plate auctions illustrates its empirical relevance.

% a bit of how the method requires few assumptions
%Our assumptions are arguably mild. The baseline version of the model assumes that the bidders are symmetric and have independent private values (IPV). We require a large number of bidders and that the valuation distribution falls within the domain of attraction of the extreme value (EV) distribution, a condition satisfied by commonly used distributions like Pareto, Student-t, Gaussian, and uniform distributions (see Section \ref{sec:setup} for details). In this setup, we can conduct inference using transaction prices from just three or more auctions. Furthermore, the auctions do not need to have the same number of bidders (which may remain unknown), but should diverge at the same rate. Beyond the IPV setup, we discuss how to extend our analysis to allow for conditional IPV and reserve price. 

% what we do in our new framework
In our benchmark model, we consider second-price sealed-bid auctions with a large number of symmetric bidders holding independent private values (IPV). 
%We first characterize the limiting behavior of the transaction prices (as the number of bidders diverges) by an extreme value distribution with a single parameter, known as the {\it tail index}. This asymptotic representation allows us to conduct asymptotically valid inference for several auction fundamentals, building on recent developments in \cite{elliott2015}.
Unlike traditional frameworks that rely on a growing number of auctions, transaction price data alone does not fully identify the valuation distribution. However, as the number of bidders diverges, the transaction price data allows us to characterize the tail properties of the bid distribution. In turn, these can be used to conduct inferences on auction fundamentals, such as the winner's expected utility and the seller's expected revenue. Additionally, analyzing tail behavior allows us to assess whether the valuation distribution has bounded support and positive density at its upper limit—common assumptions in auction theory \citep[e.g.,][]{MaskinRiley1984,guerre2000,GuerreLuo2022}. The proposed inference methods in our framework control size and demonstrate optimal power properties. Specifically, we use results in \cite{elliott2015} to guarantee that our confidence intervals minimize expected length and our hypothesis tests maximize weighted average power. 

% extension to first-price auction
Beyond the benchmark model, we consider three extensions. 
First, we show that our analysis for second-price auctions extends to first-price auctions. 
% we extend our analysis to first-price auctions and provide new insights. 
% extension to binding reserve price
Second, we broaden our analysis to encompass potential bidders and conditional IPV scenarios. The presence of potential bidders often occurs when some buyers' bids remain unsubmitted or unobserved. This scenario complicates standard methods, particularly when the count of potential bidders is unknown. In contrast, our analysis remains unaffected, as the tail of the underlying valuation distribution and the equilibrium bidding function remain unchanged. % Furthermore, we explore the impact of reserve prices on the seller's expected revenue, uncovering findings divergent from the classical understanding. Refer to Section \ref{sec:extension} for detailed insights.
% extension to unobserved heterogeneity
Third, our analysis extends to examining cases involving affiliated values. Specifically, we accommodate auction-level unobserved heterogeneity $A$, conditional on which buyers possess independent values. This scenario, known as conditional IPV, has been explored by \citet{li2000cipv}. With only a limited number of auctions, existing methods struggle to identify the underlying valuation distribution $F_{V|A}$ conditional on this heterogeneity. However, our methodology can be employed to discern the tail of $F_{V|A}$ provided we observe at least three bids from a single auction. 


% % simulation
% We emphasize that our many-bidder framework does not literally require an infinite number of bidders. Instead, it is designed to capture empirical scenarios where the number of bidders is relatively large compared to the number of auctions. We evaluate the finite sample performance of our proposed method through extensive simulations in Section \ref{sec:simulation}. Our findings indicate that the proposed confidence intervals exhibit excellent coverage and length properties as long as there are more than 10 bidders. %Additionally, the method performs well even when the number of bidders varies slightly across auctions.

% % more about the application
% We apply our proposed method to Hong Kong vehicle license plate auctions in Section \ref{sec:application}.
% First, we analyze the single-letter plate auctions and construct confidence intervals for the winner's expected utility and the government's expected revenue. Second, we examine auctions for plates with two letters and two numbers, such as {\texttt{UU28}, \texttt{UU38}, \texttt{UU68}}. While the expected utility and revenue are also high, they are significantly lower than those for the single-letter plates. Lastly, we analyze ordinary plates with two distinct letters and four distinct numbers, where the expected utility and revenue are in the range of several thousand HKD. Our test also significantly rejects the regularity conditions commonly assumed in existing nonparametric methods.

% overview
The rest of the paper is organized as follows. The following subsection reviews the existing literature and discusses our contribution. Section \ref{sec:setup} establishes the new asymptotic framework with many bidders and discusses its relationship with extreme value (EV) theory. Section \ref{sec:sp} focuses on second-price auctions and introduces our new inference methods. Section \ref{sec:sp-supplus} presents confidence intervals for the winner's expected utility, Section \ref{sec:sp-revenue} introduces confidence intervals for the seller's expected revenue, and Section \ref{sec:sp-index} outlines hypothesis tests for the tail index. Section \ref{sec:extension} extends the analysis to first-price auctions, the conditional IPV model, and the presence of reserve prices. We provide Monte Carlo simulations in Section \ref{sec:simulation}. Section \ref{sec:application} applies our inference methods to data from Hong Kong license plate auctions,
% First, we analyze the single-letter plate auctions and construct confidence intervals for the winner's expected utility and the government's expected revenue. Second, we examine auctions for plates with two letters and two numbers, such as {\texttt{UU28}, \texttt{UU38}, \texttt{UU68}}. While the expected utility and revenue are also high, they are significantly lower than those for the single-letter plates. Lastly, we analyze ordinary plates with two distinct letters and four distinct numbers, where the expected utility and revenue are in the range of several thousand HKD. Our test also significantly rejects the regularity conditions commonly assumed in existing nonparametric methods.
and Section \ref{sec:conclusion} concludes. The paper's appendix provides all proofs, auxiliary results, and computational details.

\subsection{Related literature}\label{sec:literature}
% start literature review

%Although our framework is distinct from the existing one, our work builds upon and relates to the extensive literature about auctions. We discuss our contributions and relations with the existing literature in the following several aspects and refer readers to the comprehensive literature review by \citet{gentry2018review}.

As mentioned earlier, our paper considers auctions in an asymptotic framework where the number of bidders diverges. This framework has been extensively employed in the economic theory literature \citep[e.g.,][]{hongshum2004,virag2013,ditillio2021} but less so in econometrics. \citet{HongPaarschXu2014} studies the asymptotic distribution of the transaction price in a clock model of a multi-unit, oral, ascending-price auction as the numbers of bidders and units increase. \citet{krasnokutskaya2022estimating} considers a latent group structure on the set of agents and allows both the number of agents and the number of markets to grow. \citet{cassola20132007} studies multi-object auctions and proposes within-auction bootstrap methods, which is similar to a many-bidders asymptotic analysis. \citet{menzel2013} shows that the nonparametric estimator of the valuation distribution may become irregular and perform poorly when the number of bidders increases. This issue does not affect our method, which relies on the EV approximation. %\citet{WilliamsZachariadis2021} studies the asymptotic distribution of the transaction price in Buyer's Bid Double Auction. 

In comparison, much of the existing econometric literature has focused on identifying and estimating the valuation distribution with a finite number of bidders and a diverging number of auctions. 
% In a setting with a finite number of bidders, like timber auctions, identifying the valuation distribution is the key to understanding the relevant auction features. 
The seminal work by \citet{athey2002} derives general results about identifying the valuation distribution from the distribution of bids. \citet{hailetamer2003} investigates English auctions and establishes bounds on the valuation distribution and other objects of interest with minimal structural assumptions. \cite{chesher2017wp} extends these bounds to the non-IPV setup. 
\citet{aradillas2013} nonparametrically identifies bounds on seller profit and bidder surplus while accounting for variations in the number of bidders across auctions. \citet{BrendstrupPaarsch2006} and \citet{Komarova2013} derive nonparametric identification of the valuation distribution with the transaction price and the winner's identity. \citet{BrendstrupPaarsch2007} investigates multi-object English auctions and establishes semiparametric identification using winning bids under the Archimedean copula assumption.

% This paragraph lists several papers that use at least two consecutive order statistics and the knowledge of the number of bidders. In comparison, our framework requires only one order statistic and not the number of bidders.
The existing auction literature has considered different types of data, including submitted bids, the number of actual bidders, and that of the potential bidders. 
%There is a vast literature that considers inference in the more traditional asymptotic framework in which the number of auctions diverges to infinity. We now briefly highlight some of its recent contributions. 
For example, \cite{Li2005} studies first-price auctions with entry and binding reservation prices and estimates the valuation distribution with the observed bids and the number of actual bidders. %\cite{adams2007} studies eBay auctions and requires knowing the distribution of the potential bidders. 
Without knowing the number of bidders, \cite{an2010} proposes using a proxy of the number of bidders and an instrument variable. \citet{KimLee2014}, \citet{song2015}, \citet{Mbakop2017}, and \citet{freyberger2022} construct identification of the valuation distribution with two or more order statistics of bids. \cite{shneyerov2011} derives nonparametric identification of model primitives based on finitely many groups of bidders. Recently, \citet{LuoXiao2023} derives identification results with two consecutive order statistics and an instrument or three consecutive ones. All these methods of identification, estimation, and inference are based on the traditional asymptotic framework with many auctions and multiple bids from each auction. We refer to \citet{hickman2012} and \citet{gentry2018review} for recent surveys.

% compareison with GL2022: GL2022 does not have inference
As already mentioned, our paper provides inference based only on transaction prices. This aspect of our paper resembles the recent work by \cite{GuerreLuo2022}. However, there are considerable differences between this paper and our contribution. In particular, \cite{GuerreLuo2022} considers first-price auctions in which the number of bidders in each auction is random and has finite support. In this context, they establish that the winning bid is increasing in the number of bidders and, hence, the density of the winning bids exhibits discontinuities as the number of bidders changes. Our framework has many differences with \citet{GuerreLuo2022}. First, their identification strategy relies on the assumption that the number of bidders is finite, while we specialize in the case when this amount diverges. Thus, the contributions are designed for different empirical environments. Second, their argument does not apply to second-price auctions, while ours does. Third, their method relies on observing a large number of auctions, while ours can deal with applications with a finite number of them. Finally, their paper delivers identification analysis, while ours focuses on inference.

% comparison with testing in auction models
Finally, our paper is also connected to the literature on testing in auction models. \citet{DonaldPaarsch1996} introduces parametric tests within the context of IPV setups. \citet{HaileHongShum2003} devises nonparametric tests for common values in first-price auctions. \citet{JunPinkseWan2010} develops a nonparametric test for affiliation. \citet{HortacsuKastl2012} proposes a test of common values when some bidders have information about rivals' bids. \citet{HillShneyerov2013} develops a test for common values in first-price auctions utilizing tail indices. \citet{LiuLuo2017} puts forward a nonparametric test for comparing valuation distributions in first-price auctions. 
Compared with these studies, our test in Section \ref{sec:sp-index} serves as a model specification test. Specifically, we consider testing if the support of the valuation distribution is bounded and the density is strictly positive near the upper end-point, both of which are typically imposed in the existing literature.
%Relative to these studies, our work stands out as we address testing within a comprehensive framework accommodating numerous bidders in both first and second-price auctions. Furthermore, our methods remain applicable when multiple bids from each auction are available, operating effectively after conditioning on unobserved heterogeneity and requiring only a small number of auctions.


% Application summary
% We apply our proposed methodology to the Hong Kong car license auction data from 1997 to 2008, provided by \cite{ng/chong/du:2010}. This auction is the standard ascending price auction, which is equivalent to the second-price auction under our assumptions. We interpret each license plate in our dataset as a separate instance of an auction. After restricting our dataset to homogeneous license plates (e.g., dropping special plates that contain favorable digit combinations), we obtain relatively small yearly samples. In fact, the number of homogeneous auctions per year ranges between 12 and 34, with an average of 26.5. Given these sample sizes, we cannot accurately estimate a year-specific distribution of valuations. Our confidence intervals show that there is a considerable winner's expected utility with large variability over time. This variability could be explained by the introduction of special plates introduced in 2006. Moreover, our hypothesis tests indicate that the standard regularity conditions imposed by traditional inference methods in auction literature are rejected every year in our dataset.

